\titleformat{\chapter}[display]{\normalfont\bfseries\color{BellaInk}}{\filleft\sffamily\color{BellaBlue}\fontsize{13}{16}\selectfont ANNEXE~\thechapter\enspace /\enspace 附录}{1ex}{\filleft\fontsize{27}{32}\selectfont}[\vspace{1.2ex}\titlerule]
\chapter{校勘参考资料与待核事项}
\markright{校勘参考资料与待核事项}
本附录供查阅正文校注。脚注中的“原”为原书注释，“校”为新增校勘；[E1]–[E69] 对应下列资料。仅据摘要、书目或二级资料的条目分别注明，待核线索不作为确证。

\section*{参考资料}
\begingroup
\small\setstretch{1.08}\setlength{\parindent}{0pt}\setlength{\parskip}{.38em}
\par\noindent\begin{minipage}{\linewidth}
\phantomsection\label{ref-1}\noindent\textbf{[E1] Bibliothèque nationale de France.} \textit{Blaise Pascal (1623-1662) — Bibliographie}.\par\nobreak
Pascal 生卒年的书目依据。\par
{\small\href{https://www.bnf.fr/fr/blaise-pascal-1623-1662-bibliographie}{在线资料}}\par
\par\end{minipage}\par

\par\noindent\begin{minipage}{\linewidth}
\phantomsection\label{ref-2}\noindent\textbf{[E2] University of Oxford, ETCSL.} \textit{Sumerian language}.\par\nobreak
苏美尔语与阿卡德语的分类；区分最早文字证据与语言起源。\par
{\small\href{https://etcsl.orinst.ox.ac.uk/edition2/language.php}{在线资料}}\par
\par\end{minipage}\par

\par\noindent\begin{minipage}{\linewidth}
\phantomsection\label{ref-3}\noindent\textbf{[E3] Steven Pinker.} \textit{The Language Instinct (1994/2007)}.\par\nobreak
作者书介涉及语言的先天基础与儿童习得；先天基础不表示具体语言无需学习。\par
{\small\href{https://stevenpinker.com/publications/language-instinct-19942007}{在线资料}}\par
\par\end{minipage}\par

\par\noindent\begin{minipage}{\linewidth}
\phantomsection\label{ref-4}\noindent\textbf{[E4] Jean-Jacques Hublin et al. (2017).} \textit{New fossils from Jebel Irhoud, Morocco and the pan-African origin of Homo sapiens}.\par\nobreak
Nature 546, 289–292。Jebel Irhoud 化石的年代约为 $315\pm34$ 千年前。\par
{\small\href{https://doi.org/10.1038/nature22336}{在线资料}}\par
\par\end{minipage}\par

\par\noindent\begin{minipage}{\linewidth}
\phantomsection\label{ref-5}\noindent\textbf{[E5] Cecilia S. L. Lai et al. (2001).} \textit{A forkhead-domain gene is mutated in a severe speech and language disorder}.\par\nobreak
Nature 413, 519–523。FOXP2 与言语、语言发育有关，但不足以证明单一突变产生句法。\par
{\small\href{https://doi.org/10.1038/35097076}{在线资料}}\par
\par\end{minipage}\par

\par\noindent\begin{minipage}{\linewidth}
\phantomsection\label{ref-6}\noindent\textbf{[E6] Johannes Krause et al. (2007).} \textit{The derived FOXP2 variant of modern humans was shared with Neandertals}.\par\nobreak
Current Biology 17(21), 1908–1912；PMID 17949978。研究比较了现代人与尼安德特人的特定 FOXP2 变异。\par
{\small\href{https://doi.org/10.1016/j.cub.2007.10.008}{在线资料}}\par
\par\end{minipage}\par

\par\noindent\begin{minipage}{\linewidth}
\phantomsection\label{ref-7}\noindent\textbf{[E7] Albert Mehrabian.} \textit{“Silent Messages” — Description and Ordering Information}.\par\nobreak
作者对 7\%—38\%—55\% 结果的说明：适用于情感与态度表达，不能推广为所有交流的信息比例。\par
{\small\href{https://www.kaaj.com/psych/smorder.html}{在线资料}}\par
\par\end{minipage}\par

\par\noindent\begin{minipage}{\linewidth}
\phantomsection\label{ref-8}\noindent\textbf{[E8] E. S. Savage-Rumbaugh et al. (1993).} \textit{Language comprehension in ape and child}.\par\nobreak
Monographs of the Society for Research in Child Development 58(3–4), 1–222。依据论文摘要（PMID 8366872）及 Ape Initiative 介绍，区分口语理解、图形符号与手语。\par
{\small\href{https://pubmed.ncbi.nlm.nih.gov/8366872/}{在线资料}}\par {\small\href{https://www.apeinitiative.org/kanzi}{在线资料}}\par
\par\end{minipage}\par

\par\noindent\begin{minipage}{\linewidth}
\phantomsection\label{ref-9}\noindent\textbf{[E9] NOAA, National Ocean Service.} \textit{How far does light travel in the ocean?}\par\nobreak
海洋透光深度的说明；十米不是普遍无光界限。页面更新于2024年6月16日。\par
{\small\href{https://oceanservice.noaa.gov/facts/light_travel.html}{在线资料}}\par
\par\end{minipage}\par

\par\noindent\begin{minipage}{\linewidth}
\phantomsection\label{ref-10}\noindent\textbf{[E10] Wikipedia contributors.} \textit{Gregory Bateson}.\par\nobreak
二级传记资料，仅用于生卒年初校：卒年应为1980，而非1990。\par
{\small\href{https://en.wikipedia.org/wiki/Gregory_Bateson}{在线资料}}\par
\par\end{minipage}\par

\par\noindent\begin{minipage}{\linewidth}
\phantomsection\label{ref-11}\noindent\textbf{[E11] Julien Offray de La Mettrie.} \textit{L’Homme machine}. Leyde, Élie Luzac fils, 1748.\par\nobreak
依据1748年版书名页转录中的 MDCCXLVIII；版本印年与初刊流通年须区分。\par
{\small\href{https://fr.wikisource.org/wiki/L’Homme_machine_(1748)}{在线资料}}\par
\par\end{minipage}\par

\par\noindent\begin{minipage}{\linewidth}
\phantomsection\label{ref-12}\noindent\textbf{[E12] Wiktionnaire.} \textit{discourir}, rubrique « Étymologie ».\par\nobreak
词条将词源追溯至拉丁语 discurrere，而非希腊语 logos。\par
{\small\href{https://fr.wiktionary.org/wiki/discourir}{在线资料}}\par
\par\end{minipage}\par

\par\noindent\begin{minipage}{\linewidth}
\phantomsection\label{ref-13}\noindent\textbf{[E13] Henri Estienne.} \textit{Project du livre intitulé De la precellence du langage François}. Mamert Patisson, 1579.\par\nobreak
Google Books 所收里昂公共图书馆书目记录署 Henri Estienne，而非 Robert Estienne。\par
{\small\href{https://books.google.es/books?id=5gkVrK527vYC}{在线资料}}\par
\par\end{minipage}\par

\par\noindent\begin{minipage}{\linewidth}
\phantomsection\label{ref-14}\noindent\textbf{[E14] Biblissima+.} \textit{Virgile (0070-0019 av. J.-C.)}, auteur Q16458.\par\nobreak
作者规范记录 Q16458：Virgile 生于公元前70年，供与 Denys 所列年代对照。\par
{\small\href{https://portail.biblissima.fr/fr/ark:/43093/pdataa05dbb0900f756522136441ae64be4b0d2e7de1b}{Biblissima+ — Virgile，作者规范记录 Q16458}}\par
\par\end{minipage}\par

\par\noindent\begin{minipage}{\linewidth}
\phantomsection\label{ref-15}\noindent\textbf{[E15] Académie française.} \textit{Claude Favre de Vaugelas}, notice biographique.\par\nobreak
官方传记列生卒日为1585年1月6日、1650年2月26日；原书卒年1659有误。\par
{\small\href{https://www.academie-francaise.fr/les-immortels/claude-favre-de-vaugelas}{Académie française — Claude Favre de Vaugelas}}\par
\par\end{minipage}\par

\par\noindent\begin{minipage}{\linewidth}
\phantomsection\label{ref-16}\noindent\textbf{[E16] Hachette Livre--BnF.} \textit{Lettres de Madame de Sévigné, de sa famille et de ses amis. Album}. Réimpression, 2013.\par\nobreak
出版社书目说明列 Sévigné 生卒年为1626–1696；原书1996有误。依据书目，未核书信集全文。\par
{\small\href{https://books.google.tt/books?id=9YLuoQEACAAJ}{在线资料}}\par
\par\end{minipage}\par

\par\noindent\begin{minipage}{\linewidth}
\phantomsection\label{ref-17}\noindent\textbf{[E17] National Institute of Informatics, CiNii Books.} \textit{La logique, ou, L’art de penser}, NCID BB06792048.\par\nobreak
馆藏记录列共同作者 Antoine Arnauld、Pierre Nicole；Dominique Descotes 校订，Honoré Champion, 2011，ISBN 9782745322654。\par
{\small\href{https://ci.nii.ac.jp/ncid/BB06792048}{在线资料}}\par
\par\end{minipage}\par

\par\noindent\begin{minipage}{\linewidth}
\phantomsection\label{ref-18}\noindent\textbf{[E18] Immanuel Kant.} \textit{The Critique of Pure Reason.}\par\nobreak
J. M. D. Meiklejohn 英译，Project Gutenberg 第4280号。所据为导论 I–V：先天／后验、分析／综合及先天综合的数学判断。\par
{\small\href{https://www.gutenberg.org/files/4280/4280-h/4280-h.htm}{在线资料}}\par
\par\end{minipage}\par

\par\noindent\begin{minipage}{\linewidth}
\phantomsection\label{ref-19}\noindent\textbf{[E19] Ludwig Wittgenstein.} \textit{Tractatus Logico-Philosophicus.}\par\nobreak
C. K. Ogden 英译，Project Gutenberg 第5740号德英对照本。所据条目为3.326–3.328、4.024、4.461–4.4611；Ulysses／Odysseus 例句归属仍待核。\par
{\small\href{https://www.gutenberg.org/files/5740/5740-pdf.pdf}{在线资料}}\par
\par\end{minipage}\par

\par\noindent\begin{minipage}{\linewidth}
\phantomsection\label{ref-20}\noindent\textbf{[E20] Erwin Schrödinger (1935); John D. Trimmer, trans. (1980).} \textit{The Present Situation in Quantum Mechanics.}\par\nobreak
Proceedings of the American Philosophical Society 124, 323–338。依据原论文英译第5节，区分量子叠加与经典命题的逻辑合取；链接为原网页存档。\par
{\small\href{https://web.archive.org/web/20161020061333/https://www.tuhh.de/rzt/rzt/it/QM/cat.html}{在线资料}}\par
\par\end{minipage}\par

\par\noindent\begin{minipage}{\linewidth}
\phantomsection\label{ref-21}\noindent\textbf{[E21] Bertrand Russell (1905).} \textit{On Denoting.}\par\nobreak
原论文全文，转载于 Cosma Shalizi 网站。所据为限定描述、存在与唯一性条件，以及否定辖域的分析。\par
{\small\href{https://bactra.org/Russell/denoting/}{在线资料}}\par
\par\end{minipage}\par

\par\noindent\begin{minipage}{\linewidth}
\phantomsection\label{ref-22}\noindent\textbf{[E22] Ludwig Wittgenstein; Harper Collins.} \textit{The Blue and Brown Books.}\par\nobreak
1965年版，ISBN 9780061312113。出版社书介将《蓝皮书》口授笔记定于1933–1934年；依据书介，未核全书。\par
{\small\href{https://books.google.com/books/about/The_Blue_and_Brown_Books.html?id=R7CxDSdN6mcC}{在线资料}}\par
\par\end{minipage}\par

\par\noindent\begin{minipage}{\linewidth}
\phantomsection\label{ref-23}\noindent\textbf{[E23] Katrin Amunts, Axel Schleicher, Karl Zilles (2004).} \textit{Outstanding language competence and cytoarchitecture in Broca’s speech region.}\par\nobreak
Brain and Language 89(2), 346–353；PMID 15068917。依据书目与摘要核对 Krebs 姓名；研究不支持多语能力的单一成因解释。\par
{\small\href{https://pubmed.ncbi.nlm.nih.gov/15068917/}{在线资料}}\par
\par\end{minipage}\par

\par\noindent\begin{minipage}{\linewidth}
\phantomsection\label{ref-24}\noindent\textbf{[E24] Glottolog.} \textit{Glottolog Information.}\par\nobreak
语系及孤立语的分类原则；同源关系不以亲属语言是否仍存活为定义。\par
{\small\href{https://glottolog.org/glottolog/glottologinformation}{在线资料}}\par
\par\end{minipage}\par

\par\noindent\begin{minipage}{\linewidth}
\phantomsection\label{ref-25}\noindent\textbf{[E25] Glottolog.} \textit{Family: Japonic (japo1237).}\par\nobreak
日琉语系分类；与其他语系关系未定，不表示日语没有可确证的亲属。\par
{\small\href{https://glottolog.org/resource/languoid/id/japo1237}{在线资料}}\par
\par\end{minipage}\par

\par\noindent\begin{minipage}{\linewidth}
\phantomsection\label{ref-26}\noindent\textbf{[E26] Suetonius.} \textit{Divus Iulius, §82.}\par\nobreak
The Latin Library 所收拉丁原文。tradiderunt quidam 引入的是另一传述，不是已确定的遗言。\par
{\small\href{https://www.thelatinlibrary.com/suetonius/suet.caesar.html}{在线资料}}\par
\par\end{minipage}\par

\par\noindent\begin{minipage}{\linewidth}
\phantomsection\label{ref-27}\noindent\textbf{[E27] Jonathan G. F. Powell (2007).} \textit{A new text of the Appendix Probi.}\par\nobreak
The Classical Quarterly 57(2), 687–700。待核文献线索：尚未核读校订全文，不能据此确证作者归属与年代。\par
{\small\href{https://doi.org/10.1017/S0009838807000638}{在线资料}}\par
\par\end{minipage}\par

\par\noindent\begin{minipage}{\linewidth}
\phantomsection\label{ref-28}\noindent\textbf{[E28] Académie française.} \textit{oculaire, Dictionnaire de l’Académie française, 9e édition.}\par\nobreak
依据词条检索片段：oculaire 来自晚期拉丁语 ocularis，须与 oculus 的民间音变区分。可在词典首页检索词条。\par
{\small\href{https://www.dictionnaire-academie.fr/}{在线资料}}\par
\par\end{minipage}\par

\par\noindent\begin{minipage}{\linewidth}
\phantomsection\label{ref-29}\noindent\textbf{[E29] Wikipedia contributors.} \textit{Turkish alphabet reform.}\par\nobreak
二级历史资料将改革及第1353号法律通过日期列为1928年、1928年11月1日；法律原件仍待核。\par
{\small\href{https://en.wikipedia.org/wiki/Turkish_alphabet_reform}{在线资料}}\par
\par\end{minipage}\par

\par\noindent\begin{minipage}{\linewidth}
\phantomsection\label{ref-30}\noindent\textbf{[E30] Wilhelm von Humboldt (1836).} \textit{Über die Verschiedenheit des menschlichen Sprachbaues und ihren Einfluss auf die geistige Entwickelung des Menschengeschlechts.}\par\nobreak
Dümmler, 1836。依据 Google Books 所收巴伐利亚州立图书馆书目订正1936为1836；未核该版全文。\par
{\small\href{https://books.google.com/books?id=NhBcAAAAcAAJ}{在线资料}}\par
\par\end{minipage}\par

\par\noindent\begin{minipage}{\linewidth}
\phantomsection\label{ref-31}\noindent\textbf{[E31] Matthew S. Dryer.} \textit{Order of Subject, Object and Verb.}\par\nobreak
WALS Online 第81章。所据为词汇性名词短语的取样原则、无主导语序类别及81A表；不同样本的比例不可混用。\par
{\small\href{https://wals.info/chapter/81}{在线资料}}\par
\par\end{minipage}\par

\par\noindent\begin{minipage}{\linewidth}
\phantomsection\label{ref-32}\noindent\textbf{[E32] Balthasar Bickel and Johanna Nichols.} \textit{Fusion of Selected Inflectional Formatives; Inflectional Synthesis of the Verb.}\par\nobreak
WALS Online 第20、22章。区分综合程度、语音融合与其他形态参数；同一语言不同构式可有不同的综合程度。\par
{\small\href{https://wals.info/chapter/20}{在线资料}}\par
{\small\href{https://wals.info/chapter/22}{在线资料}}\par
\par\end{minipage}\par

\par\noindent\begin{minipage}{\linewidth}
\phantomsection\label{ref-33}\noindent\textbf{[E33] J. H. Allen and J. B. Greenough; Meagan Ayer, ed.} \textit{New Latin Grammar: Uses of Prepositions, §§220–221.}\par\nobreak
Dickinson College Commentaries。in 表方向时支配宾格，表处所时支配夺格；in cubiculo sum 属后者。\par
{\small\href{https://dcc.dickinson.edu/grammar/latin/uses-prepositions}{在线资料}}\par
\par\end{minipage}\par

\par\noindent\begin{minipage}{\linewidth}
\phantomsection\label{ref-34}\noindent\textbf{[E34] ATILF–CNRTL.} \textit{hommasse, Trésor de la langue française informatisé.}\par\nobreak
hommasse 的词类、词义与派生关系；该词不是 homme 的屈折阴性形式。\par
{\small\href{https://www.cnrtl.fr/definition/hommasse}{在线资料}}\par
\par\end{minipage}\par

\par\noindent\begin{minipage}{\linewidth}
\phantomsection\label{ref-35}\noindent\textbf{[E35] IASS–AIS.} \textit{Congresses}.\par\nobreak
协会官方大会表列首届大会于1974年6月在 Milano 举行；原书1969有误。\par
{\small\href{https://iass-ais.org/congresses/}{电子核对文本}}\par
\par\end{minipage}\par

\par\noindent\begin{minipage}{\linewidth}
\phantomsection\label{ref-36}\noindent\textbf{[E36] Charles S. Peirce.} \textit{How to Make Our Ideas Clear (1878)}.\par\nobreak
原刊全文转录。所据为实践效果准则与最终共同意见的真理观，不能将个别效用直接等同于真理。\par
{\small\href{https://en.wikisource.org/wiki/Popular_Science_Monthly/Volume_12/January_1878/Illustrations_of_the_Logic_of_Science_II}{电子核对文本}}\par
\par\end{minipage}\par

\par\noindent\begin{minipage}{\linewidth}
\phantomsection\label{ref-37}\noindent\textbf{[E37] Charles S. Peirce.} \textit{A Syllabus of Certain Topics of Logic (1903)；What Is a Sign? (1894)}.\par\nobreak
Universidad de Navarra 所刊西班牙语译本。所据为 CP 2.233–272 的三分法与十类、符号定义与指示实例，以及 CP 1.180–202 的科学分类；法译引文版本另待核。\par
{\small\href{https://www.unav.es/gep/RelacionesTriadicas.html}{电子核对文本 1}}\par
{\small\href{https://www.unav.es/gep/DiversasConcepcionesLogicas.html}{电子核对文本 2}}\par
{\small\href{https://www.unav.es/gep/AnOutlineClassification.html}{电子核对文本 3}}\par
{\small\href{https://www.unav.es/gep/Signo.html}{电子核对文本 4}}\par
\par\end{minipage}\par

\par\noindent\begin{minipage}{\linewidth}
\phantomsection\label{ref-38}\noindent\textbf{[E38] Ferdinand de Saussure；Charles Bally、Albert Sechehaye 编，Albert Riedlinger 协作.} \textit{Cours de linguistique générale（1916；所查电子文本据 Payot 1971 版）}.\par\nobreak
所据为第一版序言、导论 III–V、第一部分 I–III、第二部分 IV–V，涉及三轮课程、能指／所指、任意性、线条性、关联关系与外部语言学。\par
{\small\href{https://fr.wikisource.org/wiki/Cours_de_linguistique_g%C3%A9n%C3%A9rale/Texte_entier}{电子核对文本}}\par
\par\end{minipage}\par

\par\noindent\begin{minipage}{\linewidth}
\phantomsection\label{ref-39}\noindent\textbf{[E39] Ferdinand de Saussure.} \textit{Recueil des publications scientifiques de Ferdinand de Saussure (1922)}.\par\nobreak
依据论文集书目与目录确认 Saussure 另有多篇科学论文；未逐篇核读全文。\par
{\small\href{https://fr.wikisource.org/wiki/Recueil_des_publications_scientifiques_de_Ferdinand_de_Saussure}{电子核对文本}}\par
\par\end{minipage}\par

\par\noindent\begin{minipage}{\linewidth}
\phantomsection\label{ref-40}\noindent\textbf{[E40] Harvard University, Faculty of Arts and Sciences.} \textit{Hilary Whitehall Putnam, Memorial Minute (7 November 2023)}.\par\nobreak
纪念文第一页列生卒日为1926年7月31日、2016年3月13日。\par
{\small\href{https://officeofthesecretary.fas.harvard.edu/sites/g/files/omnuum2291/files/2024-11/Putnam%20Memorial%20Minute_2023_11_07.pdf}{电子核对文本}}\par
\par\end{minipage}\par

\par\noindent\begin{minipage}{\linewidth}
\phantomsection\label{ref-41}\noindent\textbf{[E41] International Phonetic Association.} \textit{The International Phonetic Alphabet — official IPA Kiel chart (2026; revised to 2015/2005)}.\par\nobreak
依据官方网页及单页图表，辨析 \bipa{ʁ / ʀ}、\bipa{l}、\bipa{j / ɥ / w}、\bipa{ø}、鼻化符号及非肺部气流辅音；法语地区音位数不由此表确定。\par
{\small\href{https://www.internationalphoneticassociation.org/content/ipa-chart}{International Phonetic Association · 官方图表说明}}\par
{\small\href{https://www.internationalphoneticassociation.org/IPAcharts/common_files/pdfs/pdfs_IPA_charts_E/IPA_Kiel.pdf}{International Phonetic Association · 官方 IPA Kiel 图表}}\par
\par\end{minipage}\par

\par\noindent\begin{minipage}{\linewidth}
\phantomsection\label{ref-42}\noindent\textbf{[E42] J. Gordon Betts et al.; OpenStax, Rice University.} \textit{Anatomy and Physiology 2e (2022), §23.3 The Mouth, Pharynx, and Esophagus}.\par\nobreak
所据为 The Mouth 与 The Pharynx：硬腭、软腭、小舌，以及鼻咽、口咽、喉咽的区分。\par
{\small\href{https://openstax.org/books/anatomy-and-physiology-2e/pages/23-3-the-mouth-pharynx-and-esophagus}{OpenStax · §23.3（电子核对文本）}}\par
\par\end{minipage}\par

\par\noindent\begin{minipage}{\linewidth}
\phantomsection\label{ref-43}\noindent\textbf{[E43] J. Gordon Betts et al.; OpenStax, Rice University.} \textit{Anatomy and Physiology 2e (2022), §22.1 Organs and Structures of the Respiratory System}.\par\nobreak
所据为 Pharynx、Larynx、Trachea 各节：声道组成、声带所在结构与相关解剖位置。\par
{\small\href{https://openstax.org/books/anatomy-and-physiology-2e/pages/22-1-organs-and-structures-of-the-respiratory-system}{OpenStax · §22.1（电子核对文本）}}\par
\par\end{minipage}\par

\par\noindent\begin{minipage}{\linewidth}
\phantomsection\label{ref-44}\noindent\textbf{[E44] Ian Maddieson.} \textit{Tone.}\par\nobreak
WALS Online 第13章。声调定义、边缘个案与取样限制；样本比例不等于全世界语言比例。\par
{\small\href{https://wals.info/chapter/13}{电子核对文本}}\par
\par\end{minipage}\par

\par\noindent\begin{minipage}{\linewidth}
\phantomsection\label{ref-45}\noindent\textbf{[E45] Societas Linguistica Europaea.} \textit{About Societas Linguistica Europaea — History / Presidents.}\par\nobreak
协会历史与历任会长表：André Martinet 任至1967年年会，继任者为 Björn Collinder。\par
{\small\href{https://societaslinguistica.eu/about-societas-linguistica-europaea/}{电子核对文本}}\par
\par\end{minipage}\par

\par\noindent\begin{minipage}{\linewidth}
\phantomsection\label{ref-46}\noindent\textbf{[E46] Roman Jakobson.} \textit{Linguistics and Poetics.}\par\nobreak
Language in Literature 重印节选，所核页为66–70。涉及六因素／六功能、Bühler 三功能及诗学例证；法译引文版本另待核。\par
{\small\href{https://marul.ffst.hr/~bwillems/fymob/jakobsonfunctions.pdf}{电子核对文本}}\par
\par\end{minipage}\par

\par\noindent\begin{minipage}{\linewidth}
\phantomsection\label{ref-47}\noindent\textbf{[E47] Dudenredaktion.} \textit{singen; sinken.}\par\nobreak
Duden Online 音标栏：singen 为 \bipa{[ˈzɪŋən]}，sinken 为 \bipa{[ˈzɪŋkn̩]}，不构成原书所称的元音间 /k/ 与 /g/ 对立。\par
{\small\href{https://www.duden.de/rechtschreibung/singen}{电子核对文本 1}}\par
{\small\href{https://www.duden.de/rechtschreibung/sinken}{电子核对文本 2}}\par
\par\end{minipage}\par

\par\noindent\begin{minipage}{\linewidth}
\phantomsection\label{ref-48}\noindent\textbf{[E48] Simon C. Dik; De Gruyter Mouton.} \textit{The Theory of Functional Grammar, Part 1: The Structure of the Clause.}\par\nobreak
出版社记录：2nd rev. ed., 1997，ISBN 9783110154047。1997不是该卷初版；依据书目，未核正文。\par
{\small\href{https://www.degruyterbrill.com/document/doi/10.1515/9783110218367/html}{电子核对文本}}\par
\par\end{minipage}\par

\par\noindent\begin{minipage}{\linewidth}
\phantomsection\label{ref-49}\noindent\textbf{[E49] SIL Dictionary \& Lexicography Services.} \textit{Glossary.}\par\nobreak
lemma、lexeme、lexical or lexeme form 与 inflectional variants 的词典学区分；不同理论体系的术语不可直接等同。\par
{\small\href{https://sites.google.com/sil.org/dls-course/resources/glossary}{SIL · 词典学术语表}}\par
\par\end{minipage}\par

\par\noindent\begin{minipage}{\linewidth}
\phantomsection\label{ref-50}\noindent\textbf{[E50] LEOs deutsche Grammatik.} \textit{1.3.1.3.6 Partizip Perfekt mit oder ohne ge–.}\par\nobreak
条目列不带 ge– 的类别及 reagiert、gestanden、vernommen 等例，说明德语第二分词不能一律概括为 ge–…–t。\par
{\small\href{https://dict.leo.org/grammatik/OnlineGrammar/InflectionRules/FRegeln-V/Texte/Partizip-mit-ohne-ge.html}{LEO · 第二分词中的 ge–}}\par
\par\end{minipage}\par

\par\noindent\begin{minipage}{\linewidth}
\phantomsection\label{ref-51}\noindent\textbf{[E51] Paul Peter Urone; Roger Hinrichs.} \textit{College Physics 2e.} OpenStax, Rice University, 2022.\par\nobreak
所据为第4章导言、§4.3及§24.1：Isaac Newton、newton 与 Heinrich Hertz 的拼写。\par
{\small\href{https://openstax.org/books/college-physics-2e/pages/4-introduction-to-dynamics-newtons-laws-of-motion}{第4章导言}；\href{https://openstax.org/books/college-physics-2e/pages/4-3-newtons-second-law-of-motion-concept-of-a-system}{§4.3}；\href{https://openstax.org/books/college-physics-2e/pages/24-1-maxwells-equations-electromagnetic-waves-predicted-and-observed}{§24.1}}\par
\par\end{minipage}\par

\par\noindent\begin{minipage}{\linewidth}
\phantomsection\label{ref-52}\noindent\textbf{[E52] Real Academia Española; Asociación de Academias de la Lengua Española.} \textit{Diccionario de la lengua española: abogado; aguacate; talibán.}\par\nobreak
abogado 来自拉丁语 advocātus，aguacate 来自纳瓦特尔语 ahuacatl；talibán 须区分波斯语复数形式与阿拉伯语词干。\par
{\small\href{https://dle.rae.es/abogado}{abogado}；\href{https://dle.rae.es/aguacate}{aguacate}；\href{https://dle.rae.es/talibán}{talibán}}\par
\par\end{minipage}\par

\par\noindent\begin{minipage}{\linewidth}
\phantomsection\label{ref-53}\noindent\textbf{[E53] Ministère de l’Éducation nationale.} \textit{La grammaire du français : Terminologie grammaticale.}\par\nobreak
Éduscol，213页。所据为第60页名词性关系从句、第80页非动词句、第189页非人称动词补足语。\par
{\small\href{https://eduscol.education.fr/document/1872/download}{Éduscol · 电子核对文本}}\par
\par\end{minipage}\par

\par\noindent\begin{minipage}{\linewidth}
\phantomsection\label{ref-54}\noindent\textbf{[E54] Bibliothèque nationale de France; Agence bibliographique de l’enseignement supérieur.} \textit{Picoche, Jacqueline (1928–2024).}\par\nobreak
BnF FRBNF12051493 与 IdRef 028740998 均列卒日为2024年12月4日；原书2014有误。\par
{\small\href{https://catalogue.bnf.fr/ark:/12148/cb12051493n}{BnF · 人名规范记录}；\href{https://www.idref.fr/028740998}{IdRef · 人名规范记录}}\par
\par\end{minipage}\par

\par\noindent\begin{minipage}{\linewidth}
\phantomsection\label{ref-55}\noindent\textbf{[E55] Joseph Courtés.} \textit{Analyse sémiotique du discours : de l’énoncé à l’énonciation.} Hachette, 1991.\par\nobreak
Google Livres 书目记录：Joseph Courtés，302页，ISBN 9782010169106。仅据书目核对作者、题名与版次资料。\par
{\small\href{https://books.google.fr/books?id=ont0QgAACAAJ}{书目核对记录}}\par
\par\end{minipage}\par

\par\noindent\begin{minipage}{\linewidth}
\phantomsection\label{ref-56}\noindent\textbf{[E56] Lionel Spinosa.} \textit{Sciences du langage : de Platon à Chomsky.}\par\nobreak
原书第339页写1907，第340页写1902，为出生年的内部对照；Benveniste 教席任期终年仍待核。\par
\par\end{minipage}\par

\par\noindent\begin{minipage}{\linewidth}
\phantomsection\label{ref-57}\noindent\textbf{[E57] 原书所引作品（版本待核）。} \textit{Problème de la poétique de Dostoïevski.}\par\nobreak
待核文献线索：原书题名与1929年须对照俄文初版及修订本题名页，尚不能确证改名改年。\par
\par\end{minipage}\par

\par\noindent\begin{minipage}{\linewidth}
\phantomsection\label{ref-58}\noindent\textbf{[E58] John R. Searle.} \textit{Speech Acts: An Essay in the Philosophy of Language. Cambridge University Press, 1969.}\par\nobreak
依据英文书目核对题名、作者与1969年；法译本版年仍待核。\par
{\small\href{https://books.google.com/books/about/Speech_Acts.html?id=t3_WhfknvF0C}{英文书目记录}}\par
\par\end{minipage}\par

\par\noindent\begin{minipage}{\linewidth}
\phantomsection\label{ref-59}\noindent\textbf{[E59] John R. Searle.} \textit{A Taxonomy of Illocutionary Acts. In Keith Gunderson (ed.), Language, Mind, and Knowledge, Minnesota Studies in the Philosophy of Science, VII, 1975, pp.344–369.}\par\nobreak
原刊扫描。所核范围为版权页及正文第344–347、354–359页，涉及行为与动词的分类、五类行为、适配方向和施事目的。\par
{\small\href{https://semantics.uchicago.edu/kennedy/classes/f09/semprag1/searle75a.pdf}{原刊扫描}}\par
\par\end{minipage}\par

\par\noindent\begin{minipage}{\linewidth}
\phantomsection\label{ref-60}\noindent\textbf{[E60] J. L. Austin; J. O. Urmson (ed.).} \textit{How to Do Things with Words. Oxford: Clarendon Press, 1962.}\par\nobreak
所据为题名页、编者前言、第115–117页及最后一讲：1955年讲座、1962年出版；uptake 须与取效后果区分。\par
{\small\href{https://archive.org/stream/austin-how-to-do-things-with-words/Austin,%20how%20to%20do%20things%20with%20words_djvu.txt}{馆藏扫描所附检索文本}}\par
\par\end{minipage}\par

\par\noindent\begin{minipage}{\linewidth}
\phantomsection\label{ref-61}\noindent\textbf{[E61] Jacques Derrida（待核原篇）。} \textit{Signature événement contexte. In Marges de la philosophie, Paris: Minuit, 1972.}\par\nobreak
待核原篇。对应校注提出引文归属与推论强弱问题，Derrida 的完整立场仍须结合原篇上下文审定。\par
\par\end{minipage}\par

\par\noindent\begin{minipage}{\linewidth}
\phantomsection\label{ref-62}\noindent\textbf{[E62] H. P. Grice.} \textit{Meaning. The Philosophical Review, 66(3), 1957, pp.377–388.}\par\nobreak
原刊扫描的检索文本。所据为自然意义／非自然意义、说话者意图及对意图的识别。\par
{\small\href{https://semantics.uchicago.edu/kennedy/classes/f09/semprag1/grice57.pdf}{原刊扫描}}\par
\par\end{minipage}\par

\par\noindent\begin{minipage}{\linewidth}
\phantomsection\label{ref-63}\noindent\textbf{[E63] H. P. Grice.} \textit{Logic and Conversation. In P. Cole and J. L. Morgan (eds.), Syntax and Semantics, 3, 1975, pp.41–58.}\par\nobreak
所核为1989年 Studies in the Way of Words 重印本第24–33页，涉及 therefore、数量准则及 violate、flout、clash；页码不能与1975年原刊混同。\par
{\small\href{https://semantics.uchicago.edu/kennedy/classes/f09/semprag1/grice75.pdf}{所核重印文本扫描}}\par
\par\end{minipage}\par

\par\noindent\begin{minipage}{\linewidth}
\phantomsection\label{ref-64}\noindent\textbf{[E64] Vrin：Manifeste du Cercle de Vienne et autres écrits}\par\nobreak
出版社2010年版书目区分 Hans Hahn、Rudolf Carnap 等作者与协调编订者 Antonia Soulez。\par
{\small\href{https://www.vrin.fr/livre/9782711622719/manifeste-du-cercle-de-vienne-et-autres-ecrits}{来源记录}}\par
\par\end{minipage}\par

\par\noindent\begin{minipage}{\linewidth}
\phantomsection\label{ref-65}\noindent\textbf{[E65] Bloomsbury：The Bloomsbury Companion to M. A. K. Halliday}\par\nobreak
出版社介绍列全名 Michael Alexander Kirkwood Halliday，支持订正总书目中的 Mickael。\par
{\small\href{https://www.bloomsbury.com/uk/bloomsbury-companion-to-m-a-k-halliday-9781350092501/}{来源记录}}\par
\par\end{minipage}\par

\par\noindent\begin{minipage}{\linewidth}
\phantomsection\label{ref-66}\noindent\textbf{[E66] La précellence du langage françois：作者书目记录}\par\nobreak
Google Books 馆藏记录列作者 Henri Estienne；原书所列1896年版的其他出版细节仍待核。\par
{\small\href{https://books.google.fr/books?id=HqoGAAAAQAAJ&printsec=frontcover}{来源记录}}\par
\par\end{minipage}\par

\par\noindent\begin{minipage}{\linewidth}
\phantomsection\label{ref-67}\noindent\textbf{[E67] 法国国家图书馆：Traicté de la grammaire françoise (1557)}\par\nobreak
BnF 馆藏记录：Robert Estienne 著，Colette Demaizière 校注，H. Champion, 2003。\par
{\small\href{https://catalogue.bnf.fr/ark%3A/12148/cb39007830f}{来源记录}}\par
\par\end{minipage}\par

\par\noindent\begin{minipage}{\linewidth}
\phantomsection\label{ref-68}\noindent\textbf{[E68] Basic Books：Where Mathematics Come From}\par\nobreak
出版社页面与 UC San Diego 阅读书目均列 George Lakoff、Rafael E. Nunez；Rafael 为名。平装重版日期不能替换原书所列2000年。\par
{\small\href{https://www.hachettebookgroup.com/titles/george-lakoff/where-mathematics-come-from/9780465037711/?lens=basic-books}{来源记录}}\par
\par\end{minipage}\par

\par\noindent\begin{minipage}{\linewidth}
\phantomsection\label{ref-69}\noindent\textbf{[E69] Charles A. Ferguson：Diglossia。}\par\nobreak
出版社注册的 Crossref DOI 元数据列 Word 15(2), 1959, pp.325–340；原书“435 p.”有误。依据书目元数据，未核论文全文。\par
{\small\href{https://api.crossref.org/works/10.1080/00437956.1959.11659702}{DOI注册元数据}}\par
\par\end{minipage}\par
\endgroup

\section*{待核事项}
以下疑点尚不能据现有材料作确定订正，正文保留原文或以脚注标明拟校。

\begingroup
\small\setstretch{1.06}\setlength{\parindent}{0pt}\setlength{\parskip}{.25em}
\textbf{卷首与第一章。} 致谢译者署名末字尚未辨清，保留局部图像。原书第15页 Schopenhauer 引文止于“et à laquelle la vie en commun”，缺文待核；第23页 analogie 释义中的“建构位置形式”暂拟为“未知形式”。

\textbf{第二章。} 镜像识别物种清单、蜜蜂舞型的距离界限、Tinbergen 学位论文研究对象及 Ekman 情绪分类版本待核。习题 II 重复的 Bêler 保留原样。

\textbf{第三章。} Panini 年代、Denys 任教地点、早期印刷品次序、“1451年第一本书”及1569年 Estienne 语法书归属待核。Foucault 引文的 imminent/immanent 为拟校；« se rappeler de » 的规范地位及若干“最早”“所有”的概括也需限定依据。

\textbf{第四章。} 理发师故事最初刊载年、Héraclite 生卒年、Buridan 与 Ockham 的师生关系及若干版本、引文来源待核；Ulysse 例句归属与 Carnap 引文疑似缺字尚无定论。

\textbf{第五章。} Appendix Probi 作者与年代、土耳其第1353号法律原件、illatif 词源、古代引文译本及若干数量、优先权断言待核。越南语附加符号为拟校；泰米尔语及历史印欧语等例证仍需逐例审校。原书统计应结合其年代与取样口径理解。

\textbf{第六章。} Gustave Guillaume 引文卷次与页码、Morris“修订版本”的时间关系及部分法译引文来源待核。原书第164页“Charles Peirce explique :”后未见边界明确的引文。

\textbf{第七章。} Garcia 观察与论文日期、手杖故事出处、声带频率统计口径、法语地区元音系统及谚语第10、12条音标待核。

\textbf{第八章。} structure“首次出现”的范围、1928年 Manifeste 版本、Halliday 理论名称与计算机类比、部分法译引文、“声调最多”统计及阿拉伯语 r 的例示待核。

\textbf{第九章。} “非洲语言”“阿拉伯语借词”等概括、词源、首见年份及借词数量待核；历史拼法与音标须按具体语言审定。lemme、lexème、monème、morphème 应结合各理论体系理解；习题 V 重复的 inter- 保留原样。

\textbf{第十章。} Proust 句长统计及“绝对纪录”、Tesnière 的 Comment construire une phrase? 年份与相关引文待核。习题 II.5 疑似串行。Molière 人物的故意错序、夸张与伪学术措辞属于例证，不按普通排印错误处理。

\textbf{第十一章。} 1751年“首次区分”之说、百科全书 sens 条目的作者与卷次、Port-Royal 引文边界、Mill 的 apocalypse 例证及引文版本待核。connotation、sémème、classème 等术语须区分学派用法。

\textbf{第十二章。} Benveniste 教席任期终年、Bakhtine 著作的俄文题名与版本、Derrida 引文归属及多数法译引文出处待核。中文对话“主总”为疑字；会话例句可能有不同解释。

\textbf{第十三章。} Labov、Ross、Mitford 的相关记述、双语分类史、diglossie 概念史及若干 pidgin／creole 历史断言待核。

\textbf{第十四章。} Boas、Bloomfield、Sapir、Harris、Skinner、Chomsky 各阶段理论关系及若干法译引文的版本出处待核。

\textbf{第十五章。} 若干“首次”断言、神经定位的当代解释及法译引文版本待核；语言习得与语言障碍的历史表述须结合其时代理解。

\textbf{第十六章。} 认知革命史、计算主义与联结主义的发展、原型理论与具身认知的争议，以及若干法译引文版本待核。
\endgroup
